%%%%%%%%%%%%%%%%%%%%%%%%%%%%%%%%%%%%%%%%%%%%%%%%%%%%%%%%%%%%%%%%%%%%%%
% FACILITIES AND OTHER RESOURCES -- BiVRepair R01
%
% UT Southwestern content from the clinical team, Yale (Pfaller)
% section added; layout follows the submitted DVA R01 facilities
% document (site sections with bold block leads).
%%%%%%%%%%%%%%%%%%%%%%%%%%%%%%%%%%%%%%%%%%%%%%%%%%%%%%%%%%%%%%%%%%%%%%
\documentclass[11pt]{article}
\usepackage{definitions}
\begin{document}

\textbf{\uppercase{Facilities and Other Resources -- UT Southwestern Medical Center}}

UT Southwestern, one of the premier academic medical centers in the nation, integrates pioneering biomedical research with exceptional clinical care and education. The institution's faculty includes many distinguished members, including six who have been awarded Nobel Prizes since 1985. The faculty of almost 2,800 is responsible for groundbreaking medical advances and is committed to translating science-driven research quickly to new clinical treatments. UT Southwestern physicians provide medical care in about 80 specialties to more than 105,000 hospitalized patients and close to 370,000 emergency room patients and facilitate approximately 3 million outpatient visits a year. The Medical Center has three degree-granting institutions (UT Southwestern Medical School, UT Southwestern Graduate School of Biomedical Sciences, and UT Southwestern School of Health Professions), which train nearly 3,600 medical, graduate, and health profession students, residents, and postdoctoral fellows each year. Ongoing support from federal agencies, such as the National Institutes of Health, along with foundations, individuals, and corporations, provides almost \$489.4 million per year to fund faculty research. UT Southwestern has approximately 18,800 employees and an operating budget of \$3.4 billion.

\textbf{Clinical (Children's Medical Center Dallas and the Heart Center).} Children's Medical Center Dallas is a quaternary care, academic children's hospital that is one of the largest in the country. The Division of Cardiology is comprised of more than 45 cardiologists, many of whom are active in research. The Heart Center (consisting of the Division of Pediatric Cardiology and Division of Pediatric Cardiothoracic Surgery) is ranked \#5 in the USA (2026--2027 {\it U.S. News \& World Report}). Dr. Andersen (Co-I) leads single ventricle and biventricular repair surgery at Children's Health Dallas. The hospital performs over 450 index congenital heart procedures per year and receives national and international referrals for biventricular repair surgery. Dr. Andersen expects to recruit \textasciitilde{}25 patients per year, whose \gls*{icmr} datasets will be included in the study.

\textbf{Pediatric Cardiac MRI Laboratory (PCMRI).} The PCMRI Laboratory is located within Children's Health and is part of UT Southwestern Biomedical Engineering. Dr. Hussain (Co-I of this proposal) is the Director of this laboratory and of the Cardiac Cross-Sectional Imaging of the Heart Center (Children's Health and UT Southwestern). The Cardiac Modeling team is embedded within this PCMRI Laboratory and is led by Dr. Chabiniok (PI of this proposal). Dr. Reddy (Co-I) is the Director of the catheterization laboratory and performs the invasive part of the \gls*{icmr} exams, the combined catheterization and cardiac MRI procedures.

\clearpage
\textbf{\uppercase{Facilities and Other Resources -- Yale University}}

Dr. Pfaller's laboratory is located in Dunham Laboratory, conveniently between Science Hill (basic life and physical science departments) and the School of Medicine. The nearby Malone Engineering Center houses many Department of Biomedical Engineering members, including those in the biomechanics group. The intellectual community at Yale is very strong, including in cardiac mechanics, mechanobiology, and cardiovascular medicine. In addition to Biomedical Engineering, there are large, active groups in the Yale Cardiovascular Research Center, the Inter-departmental Vascular Biology and Therapeutics Program, and the Systems Biology Institute. The environment for computational cardiovascular mechanics is particularly strong, and the cardiac growth model applied in Aim~3 was developed and is maintained in Dr. Pfaller's laboratory.

\textbf{Laboratory.} Dr. Pfaller has a 224-square-foot office in Dunham Laboratory plus an adjacent 619-square-foot graduate student office in Leet Oliver Memorial Hall. These offices include desk spaces with computers, dual-monitor workstations for model construction and data analysis, and a meeting area for team discussions and code review.

\textbf{Computer.} Dr. Pfaller's laboratory includes several computers and monitors for software development, testing, and rapid prototyping, with full network access to the \gls*{ycrc} shared resources described under Equipment.

\textbf{Open-Source Software Tools.} Dr. Pfaller's laboratory maintains and develops several large-scale software projects that are open-source and accessible to the community within the SimVascular project. SimVascular provides a complete pipeline from medical image data to patient-specific cardiovascular simulation and analysis. The multi-physics 3D finite element code svMultiPhysics hosts the laboratory's homogenized constrained mixture model of cardiac growth and remodeling, the mechanistic growth formulation that Aim~3 initializes from each patient's pre-recruitment Digital Twin. The lab also develops svZeroDSolver, a fast and highly modular solver for reduced-dimensional cardiovascular models, and svSuperEstimator, a model personalization tool with built-in uncertainty quantification.

\textbf{Office.} Yale Biomedical Engineering has a well-staffed, modern office served by a computer systems specialist, pre-award and post-award grants manager, graduate student coordinator, and administrative assistants. The department supports federal grant administration including budget management, subaward management, and effort certification.

\textbf{Clinical.} Not applicable. All clinical data are acquired at the primary site and shared with Yale as de-identified datasets under a Data Use Agreement.

\textbf{Animal.} Not applicable.

\textbf{Institutional Cores and Other Resources.} Yale provides a research environment that ensures responsible conduct of research and transparency through rigorous research-ethics guidelines, regulatory compliance, quality controls, and financial accountability. The \gls*{ycrc} operates and supports the university's shared high-performance computing clusters and research storage. The Yale Library Data Services group supports the Data Management and Sharing Plan and provides repository selection consultation.

\end{document}
